\section{A physical true-Gibbs update without fine energy estimation}

\begin{lemma}[Balanced two-coin construction]\label{lem:gibbs}
Suppose a balanced physical coin has success probability
$q=\sig(\ell+z)\in[1/4,3/4]$. For $a=e^{u-z}$ the following experiment has
output law $\operatorname{Bernoulli}(\sig(\ell+u))$: choose branch 1 with
probability $a/(1+a)$ and branch 0 otherwise; toss a $q$ coin; on branch 1
accept a success and return 1; on branch 0 accept a failure and return 0;
otherwise repeat. The acceptance probability on each attempt is at least
$1/4$.
\end{lemma}
\begin{proof}
The probabilities of accepted 1 and accepted 0 are $aq/(1+a)$ and
$(1-q)/(1+a)$. Their sum is at least $\min(q,1-q)\ge1/4$ and their
normalized ratio is $aq/(aq+1-q)=\sig(\ell+u)$.
\end{proof}
This is a two-coin/Bernoulli-race identity, established machinery
\cite{barker,bernoullirace}; the new use here is its physical interface and
composition with rare-event generation.

For a requested $0<\theta\le1/4$, set
\[
 J=4\lceil\log_2(8/\theta)\rceil,\qquad
 \beta=\theta/(16J).
\]
Cap the experiment after $J$ attempts, returning a fixed fallback bit on
exhaustion. The ideal tail is at most $(3/4)^J\le e^{-J/4}<\theta/8$.
Before each read hold for
$t_{\rm hold}\ge r_{\rm cal}^{-1}\log(1/\beta)$.
This makes the physical coin within conditional TV $\beta$ of $q$.
A finite dyadic branch selector has error at most $\beta$. The two
substitutions over at most $J$ attempts add at most $2J\beta=\theta/8$.

Let the numerical values of each of the actual fields $u,z$ be within
$\xi\le\theta/4$. The computed log shift then differs by at most $2\xi$;
since $\sig$ is $1/4$-Lipschitz, the resulting infinite-factory output
probability differs by at most $\xi/2\le\theta/8$. Acceptance remains at least
$1/4$ for every positive branch weight even with the perturbed shift.
Finite exponential/selector evaluation is included in the selector error.
Thus the entire capped output has conditional TV at most $\theta$, with
slack for the declared response/readout error allocations. Standard interval
exponentials and dyadic thresholds require bit cost polynomial in
$B_0+U+\log(J/\theta)$; their cost is not proportional to $e^{B_0+U}$.

The actual service bill per requested Gibbs update is bounded by
\begin{equation}\label{eq:gibbscost}
\begin{split}
 C_G(\theta)\le{}&
 C_{\rm context}+C_{\rm field}(\theta/4)+C_{\rm slew}+C_{\rm settle}\\
 &+J\!\left[r_{\rm cal}^{-1}\log(16J/\theta)+C_{\rm read}(\beta)
            +C_{\rm timer}(\Delta_G)\right]+C_{\rm arith}.
\end{split}
\end{equation}
The hold-timing resolution $\Delta_G$ is chosen so the actual hold is at
least the specified time; rounding upwards suffices if its excess time is
charged. This method requires no focal-state reset because its bias bound
holds uniformly conditionally on every earlier experimental history.

\begin{lemma}[Balanced log-odds measurement]\label{lem:odds}
At a balanced context, a finite number of physical readings,
$O(h^{-2}\log(1/\beta_0))$, and additive-field
calibration to $O(h)$ estimate $\ell_i(s)$ to absolute error $h$, with
conditional failure probability at most $\beta_0$. Every reading holds for
$O(r_{\rm cal}^{-1}\log(1/h))$. All costs are uniform over the previous
experimental transcript.
\end{lemma}
\begin{proof}
For $h\le1$, set $a=h/128$. Hold until each conditional mean differs from
the true $q_z$ by at most $a/2$. Martingale Hoeffding on
$M\ge2a^{-2}\log(2/\beta_0)$ observations bounds the empirical error by
$a$, including the conditional-mean bias, with probability at least
$1-\beta_0$. Clip the empirical probability to $[1/8,7/8]$. Since true
$q_z\in[1/4,3/4]$, clipping cannot increase its error. The logit derivative
on the clipping interval is at most $64/7<10$, so its error is at most
$10a<h/12$. Subtract the measured field $z$, whose actual value is known
to $h/4$, and evaluate the logarithm to $h/4$. The total is less than $h$.
For $h>1$ one may use the $h=1$ procedure. The hold and every calibration,
response, read and timer service are charged.
\end{proof}

No fine log-odds estimate is used in the Gibbs update of
Lemma~\ref{lem:gibbs}. Lemma~\ref{lem:odds} is used only when an actual
relative static probability needs a local weight-ratio evaluation.
